% Section 15: Synthesis II: one calculus, three languages | Part: 2 | Depends on: sec:fields, sec:lie, sec:ga-calculus | Status: done | Original lines: 2366-2374
%!TEX root = ../../main.tex
% =====================================================================
\section{Synthesis II: one calculus, three languages}\label{sec:synth2}
% =====================================================================

\begin{synthesis}[what this section does]
In $\R^3$ with its Euclidean metric and orientation, the 19th-century vector calculus, the calculus of differential forms and the geometric calculus
are three notations for one and the same set of operations. Table~\ref{tab:dictionary3} is the dictionary; the example and the two statements below show how to use it.
\end{synthesis}
% Section 15.1: The dictionary in $\R^3$ | Part: 2 | Depends on: sec:fields, sec:lie, sec:ga-calculus | Status: done | Original lines: 2375-2415
%!TEX root = ../../main.tex
\subsection{The dictionary in $\R^3$}

% REWRITTEN: definitions of vol and of the flux 2-form added (both are used in the table before Part III).
In the table, $\vol=\dd x\wedge\dd y\wedge\dd z$ is the \emph{volume form} of $\R^3$ (it gives volume $1$ to the unit cube), and the
\emph{flux $2$-form} of a vector field $F$ is $\iota_F\vol$, computed in \eqref{eq:iotavol}; Section~\ref{sec:flux} explains why its
integral over a surface is the flux of $F$. The Hodge star on forms is the one of Section~\ref{sec:duality}.

\begin{table}[ht]
\centering
\small
\renewcommand{\arraystretch}{1.4}
\begin{tabular}{@{}p{3.0cm}p{3.4cm}p{4.2cm}p{4.0cm}@{}}
\toprule
\textbf{Object} & \textbf{Vector calculus} & \textbf{Forms} & \textbf{Geometric algebra}\\
\midrule
gradient & $\operatorname{grad}f$ & $\dd f=(\operatorname{grad}f)^\flat$ & $\vd f$\\
curl & $\operatorname{curl}F$ & $\dd F^\flat=\iota_{\operatorname{curl}F}\vol=\hodge(\operatorname{curl}F)^\flat$ & $\vd\wedge F=\pss\,\operatorname{curl}F$\\
divergence & $\operatorname{div}F$ & $\dd(\iota_F\vol)=(\operatorname{div}F)\vol$ & $\vd\lc F=\vd\cdot F$\\
flux $2$-form & $F\cdot n\,\dd A$ & $\iota_F\vol=\hodge F^\flat$ & $F\pss=\pss F$\\
cross product & $a\times b$ & $\bigl(\hodge(a^\flat\wedge b^\flat)\bigr)^\sharp$ & $(a\wedge b)\pss^{-1}$\\
$\dd\dd=0$ & $\operatorname{curl}\operatorname{grad}=0,\ \operatorname{div}\operatorname{curl}=0$ & $\dd\dd=0$ & $\vd\wedge\vd\wedge=0$\\
all first derivatives at once & --- & $\dd$ and $\codiff$ separately & $\vd F=\vd\lc F+\vd\wedge F$\\
\bottomrule
\end{tabular}
\caption{The calculus of $\R^3$ in three languages (Euclidean metric, orientation $\pss=e_1e_2e_3$, $\pss^2=-1$; forms and multivectors identified by $\dd x^i\leftrightarrow e_i$).}
\label{tab:dictionary3}
\end{table}

\begin{proposition}[three-language identities]\label{prop:threelang}
For a vector field $F=F^ie_i$ on $\R^3$ (identified with $F^\flat=F_i\,\dd x^i$):
\begin{enumerate}[label=(\roman*)]
\item $\iota_F\vol=\hodge F^\flat\ \leftrightarrow\ F\pss$, with components $F^1\dd y\wedge\dd z+F^2\dd z\wedge\dd x+F^3\dd x\wedge\dd y\ \leftrightarrow\ F^1e_2e_3+F^2e_3e_1+F^3e_1e_2$;
\item $\dd F^\flat=\iota_{\operatorname{curl}F}\vol\ \leftrightarrow\ \vd\wedge F=\pss\operatorname{curl}F$;
\item $\dd(\iota_F\vol)=(\operatorname{div}F)\vol\ \leftrightarrow\ \vd\wedge(F\pss)=(\vd\cdot F)\,\pss$.
\end{enumerate}
\end{proposition}
\begin{proof}
(i) By Theorem~\ref{thm:hodgeGA}, $\hodge F=\rev F\pss=F\pss$ for a vector; $e_1\pss=e_2e_3$, $e_2\pss=e_3e_1$, $e_3\pss=e_1e_2$ (e.g.\ $e_2e_1e_2e_3=-e_1e_2e_2e_3=-e_1e_3=e_3e_1$). Compare with the formula \eqref{eq:iotavol} for $\iota_F\vol$ of Section~\ref{sec:interior}. (ii) $\vd\wedge F=\sum_{i<j}(\partial_iF_j-\partial_jF_i)\,e_ie_j=(\partial_2F_3-\partial_3F_2)e_2e_3+(\partial_3F_1-\partial_1F_3)e_3e_1+(\partial_1F_2-\partial_2F_1)e_1e_2$, which is $\pss\operatorname{curl}F$ by the same table of $e_i\pss$; the forms side is the computation of $\dd F^\flat$ in Section~\ref{sec:dform}.
% REWRITTEN (F5): both references now point backward (sec:interior, sec:dform) instead of forward (sec:intk, sec:stokes). (iii) Expanding, $\vd\wedge(F^1e_2e_3+F^2e_3e_1+F^3e_1e_2)=(\partial_1F^1+\partial_2F^2+\partial_3F^3)e_1e_2e_3$, since $e_1\wedge e_2e_3=e_2\wedge e_3e_1=e_3\wedge e_1e_2=\pss$.
\end{proof}

\begin{example}[all three languages on one field]\label{ex:threelang}
Let $F=(-y,\,x,\,z)$ on $\R^3$. \emph{Vector calculus:} $\operatorname{curl}F=(0,0,2)$, $\operatorname{div}F=1$.
\emph{Forms:} $F^\flat=-y\,\dd x+x\,\dd y+z\,\dd z$, so $\dd F^\flat=-\dd y\wedge\dd x+\dd x\wedge\dd y=2\,\dd x\wedge\dd y=\iota_{(0,0,2)}\vol$ \checkmark; and $\iota_F\vol=-y\,\dd y\wedge\dd z+x\,\dd z\wedge\dd x+z\,\dd x\wedge\dd y$, whose $\dd$ is $\partial_zz\,\dd z\wedge\dd x\wedge\dd y=\vol$ (the other two terms contribute $\partial_x(-y)=\partial_y(x)=0$), i.e.\ $\operatorname{div}F=1$ \checkmark.
\emph{Geometric algebra:} the single product $\vd F=\vd\lc F+\vd\wedge F=1+\pss(0,0,2)=1+2\,e_1e_2$ contains both answers at once: scalar part $\operatorname{div}F=1$, bivector part $\pss\operatorname{curl}F=2e_1e_2$.
\end{example}

% Section 15.2: One reason for three identities | Part: 2 | Depends on: sec:fields, sec:lie, sec:ga-calculus | Status: done | Original lines: 2416-2422
%!TEX root = ../../main.tex
\subsection{One reason for three identities}

\begin{synthesis}[$\dd\dd=0$, ``curl grad $=0$'' and ``div curl $=0$'']
In forms, $\dd\dd=0$ because mixed partial derivatives commute while $\dd x^j\wedge\dd x^i$ is antisymmetric (Example~\ref{ex:onefunction}). In geometric algebra, $\vd\wedge\vd\wedge A=0$ for the same reason: $\vd\wedge\vd=e^i_\sharp\wedge e^j_\sharp\,\partial_i\partial_j=0$.
Translated back to vector calculus, the $\dd\dd$ on $0$-forms is $\operatorname{curl}\operatorname{grad}=0$ and on $1$-forms is $\operatorname{div}\operatorname{curl}=0$ (use $\dd\dd=0$ together with the dictionary above). All three are one fact: \emph{the second derivative of a smooth function is symmetric}. The same fact is what fails for covariant derivatives, whose antisymmetrised second derivative is the curvature (Section~\ref{sec:synth3}).
\end{synthesis}

% Section 15.3: Maxwell's equations in three languages | Part: 2 | Depends on: sec:fields, sec:lie, sec:ga-calculus | Status: done | Original lines: 2423-2439
%!TEX root = ../../main.tex
\subsection{Maxwell's equations in three languages}

\begin{center}
\small
\renewcommand{\arraystretch}{1.3}
\begin{tabular}{@{}p{3.5cm}p{5.0cm}p{5.6cm}@{}}
\toprule
\textbf{Language} & \textbf{Equations} & \textbf{Needs a metric?}\\
\midrule
Vector calculus & $\operatorname{curl}E=-\partial_tB$, $\operatorname{div}B=0$, $\operatorname{div}D=\rho$, $\operatorname{curl}H=j+\partial_tD$ & curl and div do\\
Forms (Section~\ref{sec:phys-em}) & $\dd F=0$, $\ \dd\hodge F=J$ & only $\hodge$ (and only its conformal class)\\
Geometric algebra (spacetime algebra) & $\vd F=J$ with $F=E+\pss B$ & the geometric product does; both equations are one\\
\bottomrule
\end{tabular}
\end{center}
The first form equation $\dd F=0$ (the homogeneous equations) is metric-free and is the ``wedge half'' $\vd\wedge F=0$; the second ($\dd\hodge F=J$) is the ``contraction half'' $\vd\lc F=J$, equivalent by duality (Section~\ref{sec:duality}). The sign conventions differ with the signature convention (Section~\ref{sec:conventions}).

% Section 15 (keypoints): Keypoints | Part: 2 | Depends on: sec:fields, sec:lie, sec:ga-calculus | Status: done | Original lines: 2440-2448
%!TEX root = ../../main.tex
\begin{keypoints}[Part II]
\begin{itemize}
\item Manifold $=$ patchwork of charts; tangent and cotangent spaces carry all the algebra of Part~\ref{part:algebra} pointwise; forms are the antisymmetric covariant tensor fields.
\item $\dd$ is the unique graded-Leibniz extension of the differential of functions with $\dd\dd=0$; $\Lie_X=\dd\iota_X+\iota_X\dd$.
\item In $\R^3$ grad, curl, div are $\dd$ in degrees $0,1,2$, dressed with $\flat,\sharp,\hodge$; geometric algebra packages them as one operator $\vd=\vd\cdot+\vd\wedge$.
\end{itemize}
\end{keypoints}

